\documentclass[11pt]{amsart}
\usepackage{amsmath,amssymb,amsthm}
\usepackage[margin=1.1in]{geometry}
\usepackage{array,longtable,booktabs}
\usepackage{hyperref}

\newtheorem{theorem}{Theorem}
\newtheorem{lemma}[theorem]{Lemma}
\newtheorem{corollary}[theorem]{Corollary}
\theoremstyle{remark}
\newtheorem{remark}[theorem]{Remark}
\newtheorem{example}[theorem]{Example}

\newcommand{\Inf}{\mathrm{Inf}}
\newcommand{\Var}{\mathrm{Var}}
\newcommand{\E}{\mathbb{E}}

\title{A variance form of Han's bound on the Fourier entropy}
\author{Yitzchak Shmalo}
\address{Einstein Institute of Mathematics, The Hebrew University of Jerusalem, Jerusalem, Israel}
\date{}

\begin{document}

\begin{abstract}
Han proved that every Boolean function $f$ satisfies
\[
H(f)\le \frac{3}{\ln 2}\,I(f)+\sum_k \Inf_k(f)\log_2\frac{4}{\Inf_k(f)},
\]
where $H$ is the Fourier entropy and $I$ the total influence. We show that the logarithm may be taken relative to the
variance:
\[
H(f)\le \Bigl(2+\frac{3}{\ln 2}\Bigr)I(f)+\sum_k \Inf_k(f)\log_2\Bigl(1+\frac{\Var(f)}{\Inf_k(f)}\Bigr).
\]
The bound implies Han's up to an additive $I(f)$ and is smaller by a factor of order $k$ for the AND of $k$
variables. Written with $\log_2(1/\Inf_k(f))$ in the sum, Han's bound has the constant $2+3/\ln 2$ in front of $I(f)$; Li
and Han lowered it to the sharp value $1/\ln 2$, and neither their bound nor ours implies the other. It follows that the Fourier entropy-influence conjecture holds, with constant $2+3/\ln2+\log_2(1+1/\varepsilon)$,
for every function whose influences are all at least $\varepsilon\Var(f)$, and that $H(f)\le I(f)\,(3+3/\ln 2+H(\pi))$,
where $\pi_k=\Inf_k(f)/I(f)$ is the influence distribution. The proof combines the chain rule for the Fourier entropy
along a decision tree with a one-coordinate estimate, and an adaptive version holds along any decision tree.
\end{abstract}

\subjclass[2020]{06E30, 42C10, 94A17}
\keywords{Fourier entropy-influence conjecture, Boolean functions, Fourier entropy, total influence}

\maketitle

\section{Introduction}\label{sec:intro}

Let $f:\{-1,1\}^n\to\{-1,1\}$ have Fourier expansion $f=\sum_S\hat f(S)\chi_S$. By Parseval $\mu(S)=\hat f(S)^2$ is a
probability distribution on subsets of $[n]$, and we write
\[
H(f)=\sum_S\mu(S)\log_2\frac1{\mu(S)},\qquad \Inf_k(f)=\sum_{S\ni k}\mu(S),\qquad I(f)=\sum_k\Inf_k(f),
\]
and $\Var(f)=1-\hat f(\emptyset)^2$. The influence $\Inf_k(f)$ is also the probability that flipping $x_k$ changes
$f(x)$. The Fourier entropy-influence conjecture of Friedgut and Kalai~\cite{FK96} asserts that $H(f)\le C\,I(f)$ for
an absolute constant $C$.

Several unconditional bounds of the form $H(f)\le I(f)\cdot(\text{a logarithm})$ are known. Subadditivity of
entropy gives $H(f)\le\sum_k h_2(\Inf_k(f))$, where $h_2$ is the binary entropy; Keller, Mossel and
Schlank~\cite{KMS} deduce from it that $H(f)\le 2I(f)(1+\log n-\log I(f))$. Kelman, Kindler, Lifshitz, Minzer
and Safra~\cite{KKLMS} prove that the min-entropy is $O(I(f)\log I(f))$ for functions of constant variance, with a
matching concentration of the spectrum. Chakraborty, Kulkarni, Lokam and Saurabh~\cite{CKLS} and Wan, Wright and
Wu~\cite{WWW} bound $H(f)$ in terms of decision-tree and communication complexity measures, and O'Donnell and
Tan~\cite{OT13} show that the conjecture is preserved under suitable compositions; it holds for every linear
threshold function~\cite{ShmLTF}. Arunachalam, Chakraborty,
Kouck\'y, Saurabh and de Wolf~\cite{ACKSW} bound $H(f)$ by twice the average unambiguous parity-certificate
complexity. Han~\cite[Theorem~1]{Han} proves $H(f)=O\bigl(I(f)+\sum_k\Inf_k(f)\log(1/\Inf_k(f))\bigr)$, and the last
display of his proof gives the explicit form
\begin{equation}\label{eq:han}
H(f)\le\frac{3}{\ln 2}\,I(f)+\sum_k\Inf_k(f)\log_2\frac{4}{\Inf_k(f)}
=\Bigl(2+\frac3{\ln 2}\Bigr)I(f)+\sum_k\Inf_k(f)\log_2\frac1{\Inf_k(f)} .
\end{equation}
Li and Han~\cite{LH25} proved, by an information-theoretic argument, that the constant $2+3/\ln 2$
in~\eqref{eq:han} may be replaced by the sharp value $1/\ln 2$, and that the resulting inequality holds for every
real-valued $f$ with $\sum_S\hat f(S)^2=1$. The last term of~\eqref{eq:han} is large when many influences are small. For functions of small variance this is
forced, since every influence is at most $2\min(\Pr[f=1],\Pr[f=-1])\le\Var(f)$, and then the logarithm measures the size of $\Var(f)$ rather than how
the influence is spread. Our main result replaces $1/\Inf_k$ in the logarithm by $\Var(f)/\Inf_k$.

\begin{theorem}\label{thm:main}
For every $f:\{-1,1\}^n\to\{-1,1\}$,
\[
H(f)\le c_0\,I(f)+\sum_{k=1}^n\Inf_k(f)\log_2\Bigl(1+\frac{\Var(f)}{\Inf_k(f)}\Bigr),
\]
where $c_0=2+3/\ln2=6.3281\ldots$ and a term with $\Inf_k(f)=0$ is read as $0$.
\end{theorem}

Since $\Var(f)\le1$ and $\log_2(1+1/t)\le1+\log_2(1/t)$ for $t\le1$, Theorem~\ref{thm:main} implies~\eqref{eq:han}
with $c_0+1$ in place of $c_0$. Theorem~\ref{thm:main} and the bound of Li and Han~\cite{LH25} do not imply each
other: for majority their smaller constant makes theirs the smaller bound, while for $\mathrm{AND}_k$
(Example~\ref{ex:and}) the right side of theirs is $(k-1+1/\ln 2)\,I(f)$ and that of Theorem~\ref{thm:main} is below
$7.92\,I(f)$. Two consequences follow. The first is immediate from the theorem; the second also uses
$\Var(f)\le I(f)$, the Poincar\'e inequality~\cite[Chapter~2]{ODo}.

\begin{corollary}\label{cor:eps}
If $\Inf_k(f)\ge\varepsilon\Var(f)$ for every relevant coordinate $k$, then
$H(f)\le\bigl(c_0+\log_2(1+1/\varepsilon)\bigr)I(f)$.
\end{corollary}

\begin{corollary}\label{cor:pi}
Let $\pi_k=\Inf_k(f)/I(f)$ and $H(\pi)=\sum_k\pi_k\log_2(1/\pi_k)$. Then $H(f)\le\bigl(c_0+1+H(\pi)\bigr)I(f)$.
\end{corollary}

\begin{proof}
Since $\Var(f)\le I(f)$ and $\Inf_k(f)\le I(f)$,
\[
\log_2\Bigl(1+\frac{\Var(f)}{\Inf_k(f)}\Bigr)\le\log_2\frac{2I(f)}{\Inf_k(f)}=1+\log_2\frac1{\pi_k};
\]
multiply by $\Inf_k(f)$ and sum.
\end{proof}

Corollary~\ref{cor:pi} has the shape of the bound of Keller, Mossel and Schlank, without their term
$-2I(f)\log I(f)$, which helps when $I(f)>1$ and hurts when $I(f)<1$. Since $t\log_2(1+1/t)\le t+h_2(t)$, the right
side of Theorem~\ref{thm:main} never exceeds the subadditivity bound $\sum_kh_2(\Inf_k(f))$ by more than $(c_0+1)I(f)$,
and it can be smaller by a factor of order $\log(1/\Var(f))$ when the variance is small (Example~\ref{ex:and}).

The proof is short. Fixing one coordinate splits the entropy into the average entropy of the two restrictions and
a drop term $D_k$, and the drop has an exact decomposition into an erasure entropy~\cite{VW08} and a Jensen--Shannon
divergence~\cite{Shm1}. In~\cite{Shm1} the divergence part is bounded by $\frac2{\ln2}\Inf_k$ at every coordinate,
and the conjecture is shown to be equivalent to a bounded ratio of the erasure entropy to $\Inf_k$ at one coordinate of
each restriction; that equivalence gives no bound on the erasure entropy where the ratio is large, and so no bound on
$H(f)$. The new estimate is the bound on the erasure entropy in Lemma~\ref{lem:step}, which holds at every coordinate
and loses only $\log_2(1+\Var/\Inf_k)$. With it, Lemma~\ref{lem:step} bounds the drop by $\Inf_k\,(c_0+\log_2(1+\Var/\Inf_k))$ for every
Boolean function. Summing along a decision tree that queries the coordinates in a fixed order expresses $H(f)$ as
an average of such drops over restrictions, and the bound follows from the concavity of
$(t,v)\mapsto t\log_2(1+v/t)$. Querying the coordinates adaptively gives a stronger statement
(Theorem~\ref{thm:adaptive}). Section~\ref{sec:ex} works out examples and records numerical checks.

\section{The drop of one coordinate}\label{sec:drop}

Fix a coordinate $i$ and write $f_{\pm}=f_{i\to\pm1}$ for the two restrictions, functions of the other $n-1$
variables. Put
\[
g=\tfrac12(f_++f_-),\qquad h=\tfrac12(f_+-f_-),
\]
so that $f=g+x_ih$, $\hat f(T)=\hat g(T)$ and $\hat f(T\cup\{i\})=\hat h(T)$ for $T\not\ni i$. Since $f$ is Boolean,
$h(x)\in\{-1,0,1\}$ and $h(x)\ne0$ exactly when $i$ is pivotal at $x$, so
\begin{equation}\label{eq:h}
\sum_T\hat h(T)^2=\E h^2=\Inf_i(f),\qquad |\hat h(\emptyset)|\le\E|h|=\Inf_i(f).
\end{equation}
Also $g$ is the average of $f$ over $x_i$, so $\Var(g)\le\Var(f)$, that is, $\sum_{T\ne\emptyset}\hat g(T)^2\le\Var(f)$.

Let $U$ have law $\mu$. The drop of $i$ is
\[
D_i(f)=H(f)-\tfrac12H(f_+)-\tfrac12H(f_-).
\]
The chain rule $H(U)=H(U\setminus\{i\})+H(\mathbf 1[i\in U]\mid U\setminus\{i\})$, together with the fact that the law of
$U\setminus\{i\}$ is the equal mixture of the spectral distributions of $f_+$ and $f_-$ (as
$\hat f_\pm(T)=\hat g(T)\pm\hat h(T)$, and $\tfrac12(\hat g+\hat h)^2+\tfrac12(\hat g-\hat h)^2=\hat g^2+\hat h^2$), gives
\[
D_i(f)=E_i(f)+J_i(f),
\]
where $E_i(f)=H(\mathbf 1[i\in U]\mid U\setminus\{i\})$ and $J_i(f)$ is the Jensen--Shannon divergence between the
spectral distributions of $f_+$ and $f_-$. In coordinates,
\[
E_i(f)=\sum_{T\not\ni i}\Phi\bigl(\hat g(T)^2,\hat h(T)^2\bigr),\qquad
\Phi(p,q)=p\log_2\frac{p+q}{p}+q\log_2\frac{p+q}{q},
\]
with $\Phi(p,0)=\Phi(0,q)=0$. The function $\Phi$ is $(p+q)\,h_2(q/(p+q))$; it is nonnegative, positively homogeneous of
degree one, concave (it is the perspective of the concave function $h_2$; directly, its Hessian is $-\frac{1}{(p+q)\ln2}\bigl(\begin{smallmatrix}q/p&-1\\-1&p/q\end{smallmatrix}\bigr)\preceq0$), and
nondecreasing in each argument. A concave function that is homogeneous of degree one is superadditive:
$\Phi(p+p',q+q')=2\Phi(\frac{p+p'}2,\frac{q+q'}2)\ge\Phi(p,q)+\Phi(p',q')$.

\begin{lemma}[{\cite[Lemma 3]{Shm1}}]\label{lem:jsd}
$J_i(f)\le\frac2{\ln2}\Inf_i(f)$.
\end{lemma}

\begin{lemma}\label{lem:step}
For every Boolean $f$ and every coordinate $i$, with $\delta=\Inf_i(f)$,
\[
E_i(f)\le\Bigl(2+\frac1{\ln 2}\Bigr)\delta+\delta\log_2\Bigl(1+\frac{\Var(f)}{\delta}\Bigr),
\qquad
D_i(f)\le c_0\,\delta+\delta\log_2\Bigl(1+\frac{\Var(f)}{\delta}\Bigr).
\]
\end{lemma}

\begin{proof}
Separate the term $T=\emptyset$ and apply superadditivity to the rest:
\[
E_i(f)\le\Phi\bigl(\hat g(\emptyset)^2,\hat h(\emptyset)^2\bigr)+\Phi(A,B),\qquad
A=\sum_{T\ne\emptyset}\hat g(T)^2\le\Var(f),\quad B=\sum_{T\ne\emptyset}\hat h(T)^2\le\delta,
\]
using~\eqref{eq:h}. For the first term, $h_2(x)\le2\sqrt{x(1-x)}$ gives $\Phi(p,q)\le2\sqrt{pq}$, so it is at most
$2|\hat g(\emptyset)||\hat h(\emptyset)|\le2\delta$ by~\eqref{eq:h} and $|\hat g(\emptyset)|\le1$. For the second,
monotonicity gives $\Phi(A,B)\le\Phi(\Var(f),\delta)$ and
\[
\Phi(V,\delta)=V\log_2\Bigl(1+\frac\delta V\Bigr)+\delta\log_2\Bigl(1+\frac V\delta\Bigr)\le\frac{\delta}{\ln2}+\delta\log_2\Bigl(1+\frac V\delta\Bigr),
\]
since $\log_2(1+y)\le y/\ln 2$. This proves the bound on $E_i$; the bound on $D_i$ adds Lemma~\ref{lem:jsd}, and
$2+1/\ln2+2/\ln2=c_0$.
\end{proof}

The proof uses only one distinction among the sets $T$, whether $T$ is empty. Up to the term $(2+1/\ln2)\delta$ it
sharpens the subadditivity estimate $E_i(f)\le h_2(\Inf_i(f))\approx\Inf_i(f)\log_2(1/\Inf_i(f))$ by replacing $1/\Inf_i$ with $\Var/\Inf_i$ inside the
logarithm.

\section{Summing along a decision tree}\label{sec:tree}

A decision tree that queries coordinates of $f$ until the restricted function is constant defines, for a uniformly
random input, a random sequence of restrictions $f=f_0,f_1,\dots,f_m$, where $f_j$ is $f$ restricted by the first $j$
answers and $f_m$ is constant. Let $i_j$ be the coordinate queried at step $j$. Since $H$ of a constant function is
$0$ and the next answer is uniform given the past, the definition of the drop gives, by induction on the depth,
\begin{equation}\label{eq:path}
H(f)=\E\sum_{j\ge1}D_{i_j}(f_{j-1}),\qquad I(f)=\E\sum_{j\ge1}\Inf_{i_j}(f_{j-1}),
\end{equation}
the second because $\Inf_k(f)=\tfrac12\Inf_k(f_+)+\tfrac12\Inf_k(f_-)$ for $k\ne i$ and $I(f)=\Inf_i(f)+\tfrac12I(f_+)+\tfrac12I(f_-)$
(compare the splitting lemma for influences in~\cite[Section 2]{GMP}).

\begin{theorem}\label{thm:adaptive}
For every Boolean $f$ and every decision tree as above,
\[
H(f)\le c_0\,I(f)+\E\sum_{j\ge1}\Inf_{i_j}(f_{j-1})\log_2\Bigl(1+\frac{\Var(f_{j-1})}{\Inf_{i_j}(f_{j-1})}\Bigr).
\]
\end{theorem}

\begin{proof}
Apply Lemma~\ref{lem:step} to each term of~\eqref{eq:path} and use the second identity in~\eqref{eq:path}.
\end{proof}

\begin{proof}[Proof of Theorem~\ref{thm:main}]
Use the tree that queries $1,2,\dots,n$ in this order, skipping nothing; once $f_{j-1}$ is constant the remaining
terms vanish. Then $f_{k-1}$ is $f$ with $x_1,\dots,x_{k-1}$ fixed uniformly at random, and
\[
\E\,\Inf_k(f_{k-1})=\Inf_k(f),\qquad \E\,\Var(f_{k-1})=\Var(f)-\Var\bigl(\E[f\mid x_1,\dots,x_{k-1}]\bigr)\le\Var(f).
\]
The function $\psi(t,v)=t\log_2(1+v/t)$, with $\psi(0,v)=0$, is the perspective of the concave function
$v\mapsto\log_2(1+v)$, hence jointly concave on $[0,\infty)^2$, and it is nondecreasing in $v$. By Jensen's inequality
\[
\E\,\psi\bigl(\Inf_k(f_{k-1}),\Var(f_{k-1})\bigr)\le\psi\bigl(\Inf_k(f),\E\Var(f_{k-1})\bigr)\le\psi\bigl(\Inf_k(f),\Var(f)\bigr).
\]
Summing over $k$ and using Theorem~\ref{thm:adaptive} proves the theorem.
\end{proof}

\begin{remark}
The same argument with $\Var(f_{j-1})$ replaced by $1$, using the concavity of $t\mapsto t\log_2(1+1/t)$ in place of
that of $\psi$, gives Han's bound~\eqref{eq:han} up to one unit in the constant. Han's proof runs through the moments
$\sum_S|\hat f(S)|^{2+2\tau}$ of restricted functions, revealed one coordinate at a time, and differentiates at
$\tau=0$; the route above works with the entropy directly.
\end{remark}

\section{Examples and checks}\label{sec:ex}

\begin{example}\label{ex:and}
Let $f=\mathrm{AND}_k$ take the value $-1$ only at the all-ones input. Then $\hat f(\emptyset)=1-2^{1-k}$ and
$\hat f(S)=-2^{1-k}$ for $S\ne\emptyset$, so $\Inf_i(f)=2^{1-k}$, $I(f)=k2^{1-k}$ and
$\Var(f)=2^{2-k}(1-2^{-k})$. The ratio $\Var(f)/\Inf_i(f)$ is below $2$, and Theorem~\ref{thm:main} gives
$H(f)\le(c_0+\log_23)\,I(f)<7.92\,I(f)$, while the right side of~\eqref{eq:han} is $(k+1+3/\ln2)\,I(f)$. The true value is
$H(f)/I(f)=4(k-1)/k+2/(k\ln2)+o(1/k)$, which tends to $4$. The same holds for $\mathrm{OR}_k$.
\end{example}

\begin{example}
For majority on $n$ variables, $\Var=1$ and every influence is of order $n^{-1/2}$, so Theorem~\ref{thm:main} and
Han's bound both give order $\sqrt n\log n$, while $H=O(\sqrt n)$~\cite{OWZ11}; the same happens for balanced tribes.
For these functions the variance is bounded below and the loss is in the logarithm itself. A heuristic computation
along the tree that queries the coordinates of majority in order suggests that the adaptive form does not remove it.
A proof of the conjecture along these lines needs a step estimate $D_i\le C\,\Inf_i$ at some coordinate of each
restriction, which is the one-coordinate form of the conjecture~\cite{GMP,Shm1}.
\end{example}

\textit{Numerical checks.} On $1720$ random Boolean functions on $2$ to $10$ variables, with values $-1$ taken
independently with probability $0.5$, $0.2$, $0.05$ or $0.01$, the ratio of $H(f)$ to the right side of
Theorem~\ref{thm:main} is at most $0.486$ on each of these functions. Over all nonconstant monotone Boolean functions of six variables the maximum of this ratio
is $0.4735$, and over those of seven variables it is $0.4812$, attained by $\mathrm{OR}_n$ and $\mathrm{AND}_n$; the
maxima of the ratio of $H(f)$ to the right side of~\eqref{eq:han} are $0.3870$ and $0.3935$. Over $4{,}767$ coordinate steps of random functions on $2$ to $8$
variables, with $D_i$, $E_i$ and $J_i$ computed separately, the identity $D_i=E_i+J_i$ holds to within $3\cdot10^{-15}$,
and the three bounds of Lemmas~\ref{lem:jsd} and~\ref{lem:step} are attained to within the factors $0.869$, $0.599$
and $0.677$; along $928$ random adaptive trees the ratio of $H(f)$ to the right side of Theorem~\ref{thm:adaptive} is
at most $0.484$. The scripts are in the repository listed below.

\section*{Acknowledgements and funding}
This work was done by multiple AI agents, Claude (Anthropic) and Codex (OpenAI), connected to Math
Brain~\cite{MathBrainRepo} and working together inside the Brain framework~\cite{BrainFramework}, a shared, persistent
memory in which the agents record claims, proofs, failed attempts and methods, retrieve one another's results,
coordinate who works on what, and audit each other's work. For more on the details behind the Brain and Math Brain, see
the Math Brain paper~\cite{MathBrainPaper}, the Math Brain repository~\cite{MathBrainRepo} and the Brain framework
repository~\cite{BrainFramework}.

The research presented in this paper was supported by the European Research Council (ERC) under the European Union's
Horizon Europe research and innovation programme (grant agreement No.~101041711), by the Simons Foundation as part of
the Collaboration on the Mathematical and Scientific Foundations of Deep Learning, by Heights Labs, by Convex Nexus
Capital, by the Israel Science Foundation (grant number 2258/19), and by the Israel Science Foundation (ISF Grant
4101/25). Computational resources were provided in part by Google Cloud through its education credits programme (credit
530550103).

\section*{Declaration of competing interest}
The author holds equity in Heights Labs and in Convex Nexus Capital and has a management role at Convex Nexus Capital.
These entities are included in the funding declaration above.

\section*{Declaration of generative AI and AI-assisted technologies}
The paper comes out of Math Brain, an ongoing project on the use of AI systems for research mathematics, and its
results were developed with the Brain~\cite{BrainFramework} and Math Brain~\cite{MathBrainRepo}. In November 2026 a
Claude Fable agent (Anthropic), connected to the Brain and working alongside Codex agents running OpenAI's GPT-6.1 on
the same problem, derived the decomposition of the drop, the path formula and the one-coordinate estimate of
Lemma~\ref{lem:step}; a Codex agent refereed the estimate and found the shorter bound for the empty set. Claude Opus 5.5
(Anthropic) searched the literature, proved Theorems~\ref{thm:main} and~\ref{thm:adaptive} and the corollaries, wrote
the random-function check, prepared the manuscript and audited it; a second Claude Opus 5.5 agent ran the exhaustive
computation over monotone functions on six and seven variables. These checks are not peer review. The author directed the work
and is responsible for the content of the paper and for decisions concerning its dissemination.

\section*{Code and data availability}
Code and outputs: \url{https://github.com/yspennstate/fei-variance-han-bound}.

The paper is part of the Math Brain project, whose public repository is
\url{https://github.com/yspennstate/math-brain-public}~\cite{MathBrainRepo}.

\appendix
\section{Notation}\label{app:notation}

Symbols are listed roughly in the order in which they first appear, and the last column gives that place. The letter $k$
indexes coordinates, and in $\mathrm{AND}_k$ it is the number of variables.

\begin{longtable}{@{}>{\raggedright\arraybackslash}p{0.19\textwidth}>{\raggedright\arraybackslash}p{0.54\textwidth}>{\raggedright\arraybackslash}p{0.2\textwidth}@{}}
\toprule
Symbol & Meaning & First appears \\
\midrule
\endhead
$f$ & a Boolean function $\{-1,1\}^n\to\{-1,1\}$ & Section~\ref{sec:intro} \\
$\hat f(S)$, $\chi_S$ & the Fourier coefficient of $f$ at $S\subseteq[n]$, and the character $\prod_{i\in S}x_i$ & Section~\ref{sec:intro} \\
$\mu$ & the spectral distribution, $\mu(S)=\hat f(S)^2$ & Section~\ref{sec:intro} \\
$H(f)$ & the Fourier entropy $\sum_S\mu(S)\log_2(1/\mu(S))$ & Section~\ref{sec:intro} \\
$\Inf_k(f)$ & the influence of coordinate $k$, $\sum_{S\ni k}\mu(S)$ & Section~\ref{sec:intro} \\
$I(f)$ & the total influence $\sum_k\Inf_k(f)$ & Section~\ref{sec:intro} \\
$\Var(f)$ & the variance $1-\hat f(\emptyset)^2$ & Section~\ref{sec:intro} \\
$C$ & an absolute constant, as in $H(f)\le C\,I(f)$ & Section~\ref{sec:intro} \\
$h_2$ & the binary entropy $h_2(p)=-p\log_2p-(1-p)\log_2(1-p)$ & Section~\ref{sec:intro} \\
$c_0$ & the constant $2+3/\ln2=6.3281\ldots$ & Theorem~\ref{thm:main} \\
$\mathrm{AND}_k$, $\mathrm{OR}_k$ & the AND and the OR of $k$ variables, with value $-1$ when true & Section~\ref{sec:intro} \\
$\varepsilon$ & a lower bound on the influences relative to $\Var(f)$ & Corollary~\ref{cor:eps} \\
$\tau$ & the exponent in Han's moments $\sum_S|\hat f(S)|^{2+2\tau}$ & Section~\ref{sec:tree} \\
$\pi_k$, $H(\pi)$ & the influence distribution $\Inf_k(f)/I(f)$, and its entropy $\sum_k\pi_k\log_2(1/\pi_k)$ & Corollary~\ref{cor:pi} \\
$D_i(f)$ & the drop $H(f)-\tfrac12H(f_+)-\tfrac12H(f_-)$ of coordinate $i$ (written $D_k$ in Section~\ref{sec:intro}) & Section~\ref{sec:intro} \\
$f_\pm$ & the restrictions $f_{i\to\pm1}$ & Section~\ref{sec:drop} \\
$g$, $h$ & $\tfrac12(f_++f_-)$ and $\tfrac12(f_+-f_-)$, so that $f=g+x_ih$ & Section~\ref{sec:drop} \\
$U$ & a random set with law $\mu$ & Section~\ref{sec:drop} \\
$E_i(f)$ & the erasure entropy $H(\mathbf 1[i\in U]\mid U\setminus\{i\})$ & Section~\ref{sec:drop} \\
$J_i(f)$ & the Jensen--Shannon divergence between the spectral distributions of $f_+$ and $f_-$ & Section~\ref{sec:drop} \\
$\Phi(p,q)$ & $p\log_2\frac{p+q}{p}+q\log_2\frac{p+q}{q}=(p+q)\,h_2\bigl(q/(p+q)\bigr)$ & Section~\ref{sec:drop} \\
$\delta$ & the influence $\Inf_i(f)$ of the coordinate in Lemma~\ref{lem:step} & Lemma~\ref{lem:step} \\
$A$, $B$, $V$ & $\sum_{T\ne\emptyset}\hat g(T)^2$, $\sum_{T\ne\emptyset}\hat h(T)^2$, and a value of the variance, in the proof of Lemma~\ref{lem:step} & Lemma~\ref{lem:step} \\
$f_j$, $i_j$, $m$ & the restriction of $f$ by the first $j$ answers of a decision tree, the coordinate queried at step $j$, and the step at which $f_m$ is constant & Section~\ref{sec:tree} \\
$\psi(t,v)$ & $t\log_2(1+v/t)$, with $\psi(0,v)=0$ & Section~\ref{sec:tree} \\
\bottomrule
\end{longtable}

\begin{thebibliography}{99}
\bibitem{FK96} E.~Friedgut, G.~Kalai, Every monotone graph property has a sharp threshold, Proc. Amer. Math. Soc.
124 (1996), 2993--3002.
\bibitem{ACKSW} S.~Arunachalam, S.~Chakraborty, M.~Kouck\'y, N.~Saurabh, R.~de~Wolf, Improved bounds on Fourier
entropy and min-entropy, ACM Trans. Comput. Theory 13 (2021), no.~4, 1--40, doi:10.1145/3470860; STACS 2020, arXiv:1809.09819.
\bibitem{CKLS} S.~Chakraborty, R.~Kulkarni, S.~V.~Lokam, N.~Saurabh, Upper bounds on Fourier entropy, Theoret.
Comput. Sci. 654 (2016), 92--112.
\bibitem{GMP} M.~J.~Gonz\'alez, P.~MacManus, M.~C.~Pereyra, Further evidence towards the Fourier
entropy-influence conjecture, arXiv:2606.00246, 2026.
\bibitem{Han} X.~Han, A new bound for the Fourier-entropy-influence conjecture, Combinatorica 45 (2025), no.~1,
Art.~4, doi:10.1007/s00493-024-00133-z; arXiv:2312.08271.
\bibitem{KKLMS} E.~Kelman, G.~Kindler, N.~Lifshitz, D.~Minzer, M.~Safra, Towards a proof of the
Fourier-entropy conjecture?, Geom. Funct. Anal. 30 (2020), 1097--1138.
\bibitem{KMS} N.~Keller, E.~Mossel, T.~Schlank, A note on the entropy/influence conjecture, Discrete Math. 312
(2012), 3364--3372.
\bibitem{LH25} P.~Li, G.~Han, Strengthening Han's Fourier entropy-influence inequality via an
information-theoretic proof, arXiv:2512.03117, 2025.
\bibitem{ODo} R.~O'Donnell, \emph{Analysis of Boolean Functions}, Cambridge University Press, 2014.
\bibitem{OT13} R.~O'Donnell, L.-Y.~Tan, A composition theorem for the Fourier entropy-influence conjecture, ICALP 2013,
Lecture Notes in Comput. Sci. 7965, 780--791, doi:10.1007/978-3-642-39206-1\_66.
\bibitem{OWZ11} R.~O'Donnell, J.~Wright, Y.~Zhou, The Fourier entropy-influence conjecture for certain classes of
Boolean functions, ICALP 2011, Lecture Notes in Comput. Sci. 6755, 330--341.
\bibitem{Shm1} Y.~Shmalo, The one-coordinate step in the Fourier entropy-influence conjecture,
\url{https://github.com/yspennstate/fourier-entropy-influence-one-coordinate}, 2026.
\bibitem{ShmLTF} Y.~Shmalo, Fourier entropy of threshold functions and their compositions,
arXiv:2610.16069, 2026.
\bibitem{VW08} S.~Verd\'u, T.~Weissman, The information lost in erasures, IEEE Trans. Inform. Theory 54 (2008),
5030--5058.
\bibitem{WWW} A.~Wan, J.~Wright, C.~Wu, Decision trees, protocols, and the Fourier entropy-influence conjecture, ITCS 2014,
arXiv:1312.3003.
\bibitem{MathBrainRepo} Y.~Shmalo, \emph{Math Brain: public snapshot, code and tests}, 2026,
\url{https://github.com/yspennstate/math-brain-public}.
\bibitem{MathBrainPaper} Y.~Shmalo, \emph{Math Brain: shared, auditable memory for multi-agent mathematical research},
October 2026, \url{https://github.com/yspennstate/math-brain-public/tree/main/papers/math-brain-2026-10-24}.
\bibitem{BrainFramework} Y.~Shmalo, \emph{Brain Framework}, public software repository, 2026,
\url{https://github.com/yspennstate/brain-framework-public}.
\end{thebibliography}

\end{document}
